\section{Introduction / DBMS Theory}

\subsection*{Basic Concepts}
\emph{Data}: values representing real-world facts; their meaning depends on semantics.\\
\emph{Database}: an organized collection of related data about a particular domain.\\
\emph{DBMS}: programs that store, access and manage data; an interface between users/applications and the database.\\
\emph{Data model}: concepts describing data, relationships, semantics and constraints.

\subsection*{Why Use a DBMS?}
Drawbacks of file-based systems:\\
\begin{enuminline}
\item \emph{Difficult access}: retrieving data may require custom programs or scanning entire files.
\item \emph{Data isolation}: data scattered across files in different formats are difficult to combine.
\item \emph{Integrity problems}: enforcing consistency constraints separately in every application is difficult.
\item \emph{Redundancy / inconsistency}: duplicate data waste storage; updating only some copies causes disagreement.
\item \emph{Concurrent access}: simultaneous updates can interfere, e.g.\ a lost update.
\item \emph{Atomicity problems}: a failure may leave an operation partly completed, e.g.\ money deducted but not deposited. Atomicity requires all or nothing.
\end{enuminline}

\subsection*{Abstraction Levels}
\begin{enuminline}
\item \emph{Physical}: how data are stored; files, indexes, pointers, compression.
\item \emph{Logical}: what data are stored and how they are related, independent of storage details.
\item \emph{View}: only the part of the database relevant to a user; different users may have different views.
\end{enuminline}

\subsection*{Data Independence, Schema and Instance}
\emph{Physical data independence}: changes to physical storage need not change the logical schema or applications.\\
\emph{Schema}: the database structure.\\
\emph{Instance}: actual contents at a particular time.

\subsection*{Database Languages}
\emph{DDL} (Data Definition Language): defines or changes schema, e.g.\ CREATE, ALTER, DROP.\\
\emph{DML} (Data Manipulation Language): retrieves or modifies data, e.g.\ SELECT, INSERT, UPDATE, DELETE.

\subsection*{System Architecture}
\emph{Query processor}\\
\begin{enuminline}
\item \emph{DDL interpreter}: interprets schema definitions and updates the data dictionary.
\item \emph{DML compiler}: translates queries into evaluation plans; optimization chooses a low-cost plan.
\item \emph{Query evaluation engine}: executes the plan's low-level instructions.
\end{enuminline}

\emph{Storage manager}: interface between low-level stored data and applications/queries.\\
\begin{enuminline}
\item \emph{File manager}: allocates disk space and organizes stored data structures.
\item \emph{Transaction manager}: coordinates concurrent transactions and handles failures to preserve consistency.
\item \emph{Authorization / integrity manager}: checks access rights and integrity constraints.
\item \emph{Buffer manager}: fetches disk data into memory and decides what to cache.
\end{enuminline}

\emph{Disk storage}\\
\begin{enuminline}
\item \emph{Data files}: actual database data.
\item \emph{Data dictionary}: metadata, e.g.\ tables, attributes and constraints.
\item \emph{Indexes}: structures for fast record access.
\end{enuminline}
